\SourceLeaf{440}%
\[
\begin{aligned}
d\varphi={}&\frac{Z}{1-ZZ}(YdX-XdY)\\
&+\frac{d\xi}{1-ZZ}\{\zeta Y-\eta XXZ+\xi XYZ\}\\
&-\frac{d\eta}{1-ZZ}\{\zeta X-\eta XYZ+\xi YYZ\}\\
&+d\zeta(\eta X-\xi Y).
\end{aligned}
\]
Since furthermore
\[
\begin{aligned}
\eta XXZ-\xi XYZ&=\eta XXZ+\eta YYZ+\zeta ZYZ\\
&=\eta Z(1-ZZ)+\zeta YZZ\\
\eta XYZ-\xi YYZ&=-\xi XXZ-\zeta XZZ-\xi YYZ\\
&=-\xi Z(1-ZZ)-\zeta XZZ,
\end{aligned}
\]
our entire expression becomes
\[
\begin{aligned}
d\varphi={}&\frac{Z}{1-ZZ}(YdX-XdY)\\
&+(\zeta Y-\eta Z)d\xi+(\xi Z-\zeta X)d\eta+(\eta X-\xi Y)d\zeta.
\end{aligned}
\]

\Needspace{5\baselineskip}
\begin{center}15.\end{center}
\noindent\hspace*{1.2em}The formula obtained is valid in general, whatever the curved line may be. If it is a shortest line, however, the last three terms clearly cancel, and consequently
\[
d\varphi=-\frac{Z}{1-ZZ}(XdY-YdX).
\]
It is easy to see, however, that
\[
\frac{Z}{1-ZZ}(XdY-YdX)
\]
is simply the area of the part of the surface of the auxiliary sphere bounded by the element of the line $L$, the two great circles drawn through its endpoints and $(3)$, and the element thereby cut off from the great circle through $(1)$ and $(2)$. This area is taken as positive when $L$ lies on the same side as $(3)$ with respect to $(1)(2)$ and the direction from $P$ to $P'$ is the same as that from $(2)$ to $(1)$, negative when exactly one of these conditions is reversed, and positive again when both are reversed. In other words, it is positive when one traverses the
